% Appendix B: Quick-Start Templates
\chapter{Quick-Start Templates}
\label{app:templates}

These templates are available as standalone files in the \code{templates/}
directory of the repository. Copy them into your project and customize.
Each template includes WHY comments explaining the purpose of each section.

\section{Minimal CLAUDE.md (10 lines)}

\begin{minted}{markdown}
# Project Name

## Build
# WHY: Claude cannot guess non-standard build commands
- `npm run build` — compile
- `npm test` — run tests

## Code Style
# WHY: Only include rules that differ from standard conventions
- TypeScript strict mode
- Prettier for formatting

## Testing
# WHY: Verification criteria improve every interaction
- Run tests after every code change
- Never commit without passing CI
\end{minted}

\section{Production CLAUDE.md (50 lines)}

\begin{minted}{markdown}
# Project Name

## Build Commands
# WHY: Every command Claude needs to verify its own work
- `npm run build` — compile TypeScript
- `npm test` — run Jest tests
- `npm run lint` — ESLint + Prettier check
- `npm run typecheck` — tsc --noEmit

## Architecture
# WHY: Tells Claude WHERE things go, not just what they are
- src/routes/ — API endpoints (one file per resource)
- src/services/ — Business logic (no HTTP concerns)
- src/models/ — Database models (Prisma)
- tests/ — mirrors src/ structure

## Code Conventions
# WHY: Only rules Claude wouldn't infer from the codebase
- TypeScript strict mode, no `any`
- Functional style: prefer pure functions, immutable data
- Error handling: throw typed errors, never swallow exceptions

## Testing Standards
# WHY: These lines improve every interaction (Ch 6)
- Run tests after every code change
- Never commit without passing lint + tests + typecheck
- Coverage target: 80% for src/ modules
- Every new function needs: happy path + error case + edge case

## Git Conventions
- Conventional commits: feat/fix/refactor/test/docs
- PR description includes: Summary, Test Plan, Breaking Changes

## Security
# WHY: Defense in depth — CLAUDE.md + deny rules + hooks
- Never read .env files
- Never commit secrets, API keys, or credentials
\end{minted}

\section{Settings Template}

\begin{minted}{json}
{
  "permissions": {
    "allow": [
      "Read", "Grep", "Glob",
      "Write(**)", "Edit(**)",
      "Bash(npm run *)", "Bash(npx *)",
      "Bash(git status)", "Bash(git diff *)",
      "Bash(git log *)", "Bash(git add *)"
    ],
    "deny": [
      "Read(.env*)",
      "Read(secrets/**)",
      "Read(**/*.pem)",
      "Read(**/*.key)",
      "Read(**/credentials*)",
      "Bash(rm -rf /)",
      "Bash(sudo *)"
    ],
    "ask": [
      "Bash(git commit *)",
      "Bash(git push *)"
    ]
  }
}
\end{minted}

{\footnotesize \textit{Note: Permission rules use glob syntax (\code{*}),
not colon syntax. Evaluation order: deny first, then ask, then allow.}}

\section{Hook Configuration}

\begin{minted}{json}
{
  "hooks": {
    "PostFileWrite": [{
      "type": "command",
      "command": "npx prettier --write \"$FILE\"",
      "description": "Auto-format written files"
    }],
    "PreCommit": [{
      "type": "command",
      "command": "npm run lint && npm test",
      "description": "Block commits with lint errors or test failures"
    }]
  }
}
\end{minted}

\section{Skill Template}

\begin{minted}{yaml}
---
name: deploy-check
description: Pre-deployment readiness verification
allowed-tools: [Bash, Read, Glob, Grep]
---
# Pre-Deploy Check
# WHY: Catches deployment blockers before they reach production
1. Verify all tests pass: `npm test`
2. Check for uncommitted changes: `git status`
3. Verify version bump in package.json
4. Check CHANGELOG.md has entry for this version
5. Report readiness status (READY / NOT READY)
   with specific blockers if NOT READY
\end{minted}

{\footnotesize \textit{Note: Skills use \code{allowed-tools} (hyphenated).}}

\section{Subagent Definition Template}

\begin{minted}{markdown}
# .claude/agents/code-reviewer.md
---
name: code-reviewer
description: Review code changes for quality issues
tools: [Read, Grep, Glob]
---
# WHY: Isolated context prevents review bias from accumulated session context
Review the provided code changes for:
1. Logic errors and bugs
2. Missing error handling
3. Security vulnerabilities
4. Performance concerns
5. Style violations

Report findings with severity: Critical / High / Medium / Low.
Include file:line references for each finding.
\end{minted}

{\footnotesize \textit{Note: Subagents use \code{tools} (not \code{allowed-tools} or \code{allowed\_tools}).}}

\section{Rules File with Paths}

\begin{minted}{markdown}
---
paths: backend/**/*.py
---
# Backend Python Standards
# WHY: These rules only load when editing backend Python files,
# keeping context clean when working on other parts of the project
- Use async/await for all I/O operations
- SQLAlchemy 2.0 query syntax (not legacy)
- Type hints on all function signatures
- Docstrings on all public functions
- Run `pytest backend/tests/` after changes
\end{minted}

\section{Custom Command Template}

\begin{minted}{markdown}
# .claude/commands/check.md
# WHY: Encodes a decision workflow, not just a shortcut
Run a comprehensive check of the project:

1. Run the test suite: `npm test`
2. Run the linter: `npm run lint`
3. Run the type checker: `npm run typecheck`
4. Report results as a summary table:
   | Check | Status | Details |
   |-------|--------|---------|
   | Tests | PASS/FAIL | X passed, Y failed |
   | Lint  | PASS/FAIL | N warnings |
   | Types | PASS/FAIL | N errors |
\end{minted}
